题面所说的"鱼类资源量增加"只能说明恢复已经发生，还需要回答增长由什么驱动，以及基础资源的承载压力怎样变化。若把四大家鱼并成一个消费者功能群，再将水草、浮游植物、浮游动物和底栖饵料合成一个资源总量，模型仍可贴合2025年资源量恢复至1.8倍的观测值，却看不出首先承压的是哪一层资源，对"水下荒漠"这类底层生态退化也缺少解释。基于这一缺口，后续模型保留四类基础资源及其与四大家鱼的基本食性对应关系。
